%% Answer diagnostics, deployment acceptance and historical boundary cases.
\subsection{E6: Answer diagnostics}
\label{sec:e5}
\label{sec:agentic}\label{sec:abstention}\label{sec:composition}\label{sec:eval-agent}
In the platform-only check, the system answered ten paper-corpus questions
with the five relevant collections predefined. All ten responses completed
and included the annotated paper identifier; nine contained all designated
terms. Mean wall time was 131.2~s and total recorded cost was \$3.10. The
broad marker rule counted all ten responses, but it also recognises bare
identifiers. Masking the gold identifier left a marker in eight answers;
a separate rule requiring an explicit repository label or link found three.
Neither rule tests whether a citation resolves or supports a claim. This
check used the original BQ06 prompt and had no concurrent file or chunk
comparison. Tables~S14--S15 report the original and common-library grids
as separate diagnostics.

\subsection{E7: Deployment checks}
\label{sec:sysarms}
A one-document API fixture recorded six passing assertions.
The keyword check matched the document name and the catalogue check found
its collection. The vector assertion could pass on a response with an empty
results list, so the pass flag does not establish a vector hit. The contract
suite initially passed 11 of 13 checks; the two document-read checks passed
when repeated after fixture creation. Separate answer probes used direct
reads when the verifier was unavailable; they do not retest the verifier.
These checks establish selected interface behaviour, not retrieval accuracy.
In the separate catalogue-truncation case, an incomplete preview blocked
an unpinned search and pagination recovered the affected query.

\subsection{E8: Historical candidate-ordering diagnostic}
\label{sec:e6}
\label{sec:inflation}\label{sec:eval-verify}
The archived industrial-corpus diagnostic illustrates why admission and
ordering should be evaluated separately: sorting retained candidates by
verifier score moved some gold documents outside the read window.
Figures~S4--S5 and S7 report the historical ranks, separate score archive
and limits of the replay. The deployment fixtures do not repeat this
comparison or estimate the effect of the admission threshold.

\subsection{E9: Operational boundaries}
\label{sec:e7}
In an earlier cross-library failure audit, global search skipped eight
unreadable index collections but reported success. Reindexing restored a
missing global hit in one tested library. A separate read-only inspection
checked for persisted records in the application and storage layers,
including session data, vector records, and collection metadata. Eleven
assertions passed. The inspection did not include a restart or recovery test
and did not retest the missed-library case.
Historical load sweeps and fixed-candidate and library-scope probes remain
separate (Table~S10 and the companion inventory). In particular, a request
unfinished at the end of a load-test observation window is not necessarily
an explicit timeout. The present evidence does not establish a current
capacity limit, concurrent-write safety or crash recovery.
